\documentclass[a4paper, 10pt, conference]{ieeeconf}
\IEEEoverridecommandlockouts
\overrideIEEEmargins

\usepackage[caption=false, font=footnotesize]{subfig}
\usepackage{graphicx}
\usepackage{hyperref}
\hypersetup{breaklinks=true}
\usepackage{xurl}
\usepackage{amsmath}
\usepackage{amssymb}
\usepackage{booktabs}
\usepackage{multirow}
\usepackage{array}
\usepackage{float}

% A self-contained algorithm float. Built from base LaTeX and float
% rather than algorithm2e so the document compiles on any TeX
% installation without an extra package having to be present.
\newfloat{algorithm}{htbp}{loa}
\floatname{algorithm}{Algorithm}
\newcounter{algln}
\newcommand{\algreset}{\setcounter{algln}{0}}
\newcommand{\albox}[1]{% one numbered line, #1 already indented
  \refstepcounter{algln}%
  \noindent\makebox[1.5em][r]{\scriptsize\thealgln:}\;#1\par\vspace{0.4pt}}
\newcommand{\alplain}[1]{\noindent\makebox[1.5em][r]{}\;#1\par\vspace{0.4pt}}
\newcommand{\alind}{\hspace*{1.1em}}
\newcommand{\alkw}[1]{\textbf{#1}}
\newcommand{\alcom}[1]{{\footnotesize\itshape\;// #1}}
\newenvironment{algbody}{\par\vspace{1pt}\hrule\vspace{2pt}\algreset
  \scriptsize\setlength{\parindent}{0pt}}
  {\vspace{2pt}\hrule\vspace{1pt}}
\usepackage{stfloats}
\renewcommand{\arraystretch}{0.90}

% Let full-width floats occupy more of a page before LaTeX gives up and
% defers them, and allow a little more text on a page carrying one. The
% defaults are tuned for pages with at most one small float; this paper
% has four double-column figures and seven tables, and the stock values
% leave usable space empty at the foot of several pages.
\renewcommand{\topfraction}{0.92}
\renewcommand{\bottomfraction}{0.7}
\renewcommand{\textfraction}{0.07}
\renewcommand{\floatpagefraction}{0.82}
\renewcommand{\dbltopfraction}{0.92}
\renewcommand{\dblfloatpagefraction}{0.82}
\setcounter{topnumber}{3}
\setcounter{bottomnumber}{2}
\setcounter{totalnumber}{5}
\setcounter{dbltopnumber}{3}

% Whitespace around floats and captions. The stock separations assume a
% page with one or two floats; this paper carries fifteen, so the
% defaults spend close to a page on gaps alone. Nothing below changes
% any figure, table, number or sentence.
\setlength{\textfloatsep}{6pt plus 2pt minus 2pt}
\setlength{\dbltextfloatsep}{6pt plus 2pt minus 2pt}
\setlength{\intextsep}{5pt plus 2pt minus 2pt}
\setlength{\dblfloatsep}{7pt plus 2pt minus 3pt}
\setlength{\floatsep}{5pt plus 2pt minus 2pt}
\setlength{\abovecaptionskip}{4pt}
\setlength{\belowcaptionskip}{0pt}

% Display-equation and heading whitespace. Same principle: gaps only.
\AtBeginDocument{%
  \setlength{\abovedisplayskip}{4pt plus 2pt minus 2pt}%
  \setlength{\belowdisplayskip}{4pt plus 2pt minus 2pt}%
  \setlength{\abovedisplayshortskip}{2pt plus 1pt minus 1pt}%
  \setlength{\belowdisplayshortskip}{3pt plus 1pt minus 1pt}%
}
% Character protrusion and font expansion. This sets the same words in
% slightly less space by adjusting inter-letter and inter-word fitting,
% typically 2-4% over a paper this length. No word is removed and no
% type size changes.
% tracking is off deliberately: it letterspaces the small-caps table
% captions, which leaves numbers such as 94.2 unsearchable in the PDF.
\usepackage[protrusion=true, expansion=true, tracking=false]{microtype}

% Close the gap between reference entries. All 28 stay, in order.
\usepackage{etoolbox}
\apptocmd{\thebibliography}{%
  \setlength{\itemsep}{0pt plus 0.3pt}%
  \setlength{\parsep}{0pt}%
  \setlength{\topsep}{2pt}%
}{}{}

% Fixed so the title page does not silently change on every recompile.
\newcommand{\submissiondate}{September 2026}

\usepackage{titlesec}
\titlespacing*{\section}{0pt}{5pt plus 1pt minus 1pt}{3pt}
\titlespacing*{\subsection}{0pt}{4pt plus 1pt minus 1pt}{2pt}
\titlespacing*{\subsubsection}{0pt}{3pt plus 1pt minus 1pt}{2pt}

\title{\LARGE \bf
Autonomous Decentralised Search and Coverage with\\
Unmanned Aerial Systems Swarms
}

\author{Reddy, Mahadev Devireddy Venkata}

\begin{document}

%% ===================== COVER PAGE (excluded from page count) ===
\begin{titlepage}
    \begin{center}
        \vspace*{1cm}
        {\huge Autonomous Decentralised Search and Coverage with
        Unmanned Aerial Systems Swarms}

        \vspace{0.6cm}
        {\large By}

        \vspace{0.4cm}
        \textbf{Mahadev Reddy Devireddy Venkata}

        \vspace{1.2cm}
        {\huge\textbf{MSc Robotics Dissertation}}

        \vspace{16mm}

        \begin{center}
            
\includegraphics[scale=0.5]{logos/bristolcrest_colour.pdf}
            \hspace{8mm}
            
\includegraphics[scale=0.30]{logos/UWE_insignia.png}
        \end{center}

        \vspace{10mm}
        {\large School of Engineering Mathematics and Technology\\
        \textsc{University of Bristol}}\\
        \&\\
        {\large School of Engineering\\
        \textsc{University of the West of England}}\\

        \vspace{0.6cm}
        \begin{minipage}{10cm}
            \centering
            A MSc dissertation submitted to the University of
            Bristol and the University of the West of England in
            accordance with the requirements of the degree of
            \textsc{Master of Science in Robotics} in the Faculty
            of Engineering.
        \end{minipage}\\

        \vspace{0.5cm}
        Student ID: 2823511\\
        Supervisor: Dr Davood Asadi

        \vspace{0.6cm}
        \submissiondate
    \end{center}

    \newpage

    {\large{Declaration of own work}}

    I declare that the work in this MSc dissertation was carried
    out in accordance with the requirements of the University's
    Regulations and Code of Practice for Research Degree Programmes
    and that it has not been submitted for any other academic award.
    Except where indicated by specific reference in the text, the
    work is the candidate's own work. Work done in collaboration
    with, or with the assistance of, others, is indicated as such.
    Any views expressed in the dissertation are those of the author.

    \begin{flushright}
    Mahadev Reddy Devireddy Venkata, \submissiondate
    \end{flushright}

    \vspace{0.6cm}
    {\large{Ethics statement}}

    This project did not require ethical review as determined by my
    supervisor, Dr Davood Asadi. All work was simulation-based, with
    no human participants, animal subjects or physical hardware
    trials.

    \begin{flushright}
    Mahadev Reddy Devireddy Venkata, \submissiondate
    \end{flushright}

    \vspace{0.6cm}
    {\large{AI Usage Statement}}

    %% ------------------------------------------------------------
    %% ACTION REQUIRED BEFORE SUBMISSION.
    %% Rewrite this so it describes what you actually did, and check
    %% the wording your programme handbook requires. Do not sign the
    %% template below without reading it against your own practice.
    %% ------------------------------------------------------------
    Third-party AI tools were used in this project for debugging
    assistance with Python simulation code, for clarifying concepts
    in reinforcement learning and classical control, for the code and
    data audit reported in Section~\ref{sec:validity}, and for
    structural and language feedback on written drafts. All
    experimental results were generated by code written for this
    project.

    \vspace{0.6cm}
    {\large{Acknowledgements}}

    I thank Dr Davood Asadi for supervision throughout the project,
    and in particular for the feedback that prompted the extended
    evaluation in Section~\ref{sec:extended} and the diagnostic work
    in Section~\ref{sec:validity}.

    \vspace{0.6cm}
    {\large{Supplementary material}}

    Source code, all result files, the complete statistical output and
    the full figure set accompany this report as supplementary
    material.

\end{titlepage}

\maketitle
\thispagestyle{empty}
\pagestyle{empty}

%% ===================== ABSTRACT ================================
\begin{abstract}
UAV swarms that operate without a central controller are attractive
for search and coverage in disaster response, due to the fact that
losing any single agent will not end the mission. In this work the
authors ask whether a decentralized multi-agent reinforcement learning
policy (i.e., one that executes locally and passes no messages) can
match classical reactive control on area coverage. We compared four
controllers on a $50\times50$ grid representing
$150\,\text{m}\times150\,\text{m}$ with six UAVs. These were as
follows: 1. Boids flocking; 2. Artificial Potential Fields (APF); 3. a
single-agent PPO hybrid; and 4. a parameter-shared multi-agent PPO
policy (i.e., one in which all six UAVs learn and none is hand-coded).
Obstacle density was varied from 0\% to 30\%, and every condition was
run on 50 matched seeds across two disjoint seed sets, which provides
1{,}600 runs in total. Each seed will reproduce an identical obstacle
layout for every controller. Therefore the comparisons utilize paired
Wilcoxon signed-rank tests with Holm-Bonferroni correction,
matched-pairs rank-biserial correlation, and bootstrap confidence
intervals. The decentralized MARL policy exceeds the 90\% coverage
target in all eight density-by-seed-set cells by mean, and it does so
significantly in seven of them after correction. Thus it matches APF
and outperforms the two other controllers, which never reach the
target. It beats APF significantly in open terrain (i.e.,
approximately +4.7 and +5.0 percentage points at 0\% obstacles), and
it is statistically indistinguishable from it at 10\% and 20\% on both
seed sets. However, APF will retain an advantage at the hardest
density on the development set ($-3.8$ points). All controllers
transferred to held-out layouts within 2.4 points, and none did so
significantly. In addition, auditing the apparatus produced three
further findings. 1. Changing only the reward scheme, from per-agent
to the shared team reward that the multi-agent literature warns
against, costs 30.6 points, which exceeds the gap between any two
controllers compared here. 2. Repeating training at four seeds at 20\%
density shows that the one-million-step budget is "undertrained"
there: the across-seed spread represents 7.8 points, closing to 4.2 at
1.85M steps, where every seed will overtake APF. Therefore at that
density the budget, and not the method, sets the measured parity. 3.
Absolute-cell accounting shows that APF's apparent improvement with
obstacle density represents a denominator artefact, since it covers
640 fewer cells at 30\% density than at 0\%, while the free area
shrinks faster. The same audit retracts a figure from an earlier draft
of this work whose checkpoint could not be reproduced, and it reverses
the conclusion that was drawn from it.
\end{abstract}

%% ===================== INTRODUCTION ============================
\section{INTRODUCTION}\label{sec:intro}

\subsection{Motivation}

Covering an unknown area with a team of unmanned aerial vehicles
represents a long-standing problem, and it has immediate use in
disaster response, wildfire monitoring, and post-event damage
assessment~\cite{erdelj2017}. After an earthquake or flood, searching
the affected ground quickly is what will determine whether survivors
are found. A single UAV is too slow and too fragile. A centrally
controlled team is faster. However, the controller then becomes a
"single point of failure", and in a collapsed or contested environment
the communication link that centralization depends upon is exactly
what will fail first. Therefore decentralized swarms avoid that
problem, due to the fact that each agent will decide based upon what
it senses locally. In addition, the swarm will keep operating when
agents drop out~\cite{dorigo2021}.

\subsection{Problem Statement}

A decentralized swarm must cover as much of an unknown, obstructed
area as possible within a fixed flight budget, utilizing only local
sensing and without exchanging messages. There are two families of
solution. 1. Reactive controllers (e.g., Boids and Artificial
Potential Fields) are cheap, they require no training, and they
represent the de facto baselines. 2. Learned policies will provide
better coordination in principle. However, they face a scaling
barrier, due to the fact that a joint action space over $N$ agents
with $|A|$ actions each will grow as $|A|^N$, which is intractable
beyond a handful of agents. Therefore one question remains unresolved:
whether a learned decentralized policy can actually beat a well-tuned
reactive controller on this task, under matched conditions and with
evidence strong enough to believe.

\subsection{Research Gap}

There were three gaps in the existing work that motivated this
project. 1. Most studies of decentralized coverage will perform their
testing in empty or lightly obstructed environments~\cite{zhu2023}.
Therefore it is unclear how these methods behave as the environment
hardens. 2. Almost all of them will report coverage percentage and
then stop, even though battery use, collision risk, and computational
cost matter just as much once a system leaves simulation. 3. Least
discussed of all, the evidentiary standard is low. Of the eight
comparable studies in Section~\ref{sec:related}, none reports a
significance test, none reports an effect size, and none reports more
than ten replications per condition. A learned method that is reported
as beating a baseline by two points, based upon a single run, against
a classical controller whose parameters are not stated, is not a
result that a reader can act upon.

\subsection{Research Questions}

\textbf{RQ1.} Can a decentralized multi-agent PPO policy (i.e., one in
which every UAV learns from local observations and sends no messages)
reach the 90\% coverage criterion across obstacle densities from 0\%
to 30\%?

\textbf{RQ2.} How will it compare against tuned Boids and APF
controllers under matched obstacle layouts, and will the comparison
survive correction for multiple testing?

\textbf{RQ3.} Will the controllers transfer to obstacle layouts that
were unseen during development? In addition, how will they compare on
the secondary measures that determine deployability (i.e., energy,
collision risk, computational cost, load balance, and swarm-size
scaling)?

\subsection{Contributions}

\begin{enumerate}
    \item A four-way comparison of Boids, APF, a single-agent PPO
    hybrid, and a parameter-shared multi-agent PPO policy. All four
    will run on identical matched obstacle layouts, with 50
    replications per condition, paired significance tests, effect
    sizes, and bootstrap intervals. None of the comparable studies
    surveyed will report the last three.
    \item A working decentralized MARL coverage system that matches or
    beats a tuned APF baseline. There are two design decisions that
    make it learnable: 1. a per-agent reward that repairs credit
    assignment; and 2. a rollout length short enough to provide PPO
    with sufficient gradient updates.
    \item A direct measurement of the credit-assignment problem that
    the multi-agent literature describes. The authors held network,
    budget, seed, and dynamics fixed, and changed only the reward
    scheme. That change alone costs 30.6 percentage points of
    coverage, which is larger than the gap between any two
    controllers compared here.
    \item Reporting of variance across training runs, which the
    surveyed literature omits. In addition, we provide evidence that
    both the variance and the parity with APF at intermediate density
    represent artefacts of the training budget.
    \item Evidence that percentage coverage over free cells is a
    misleading metric when free area varies with obstacle density, and
    absolute-cell accounting that will correct it.
    \item A validation audit of the apparatus itself, with three
    repaired defects, one retracted figure, and one withdrawn
    experiment documented in Section~\ref{sec:validity}.
\end{enumerate}

%% ===================== RELATED WORK ============================
\section{RELATED WORK}\label{sec:related}

\textbf{Bio-inspired flocking.} Reynolds~\cite{reynolds1987}
introduced Boids as a distributed model in which each agent will
follow local rules that are computed from neighbors inside a
perception radius, with no global state and no central controller. It
represents the origin of the claim that coherent group behavior will
emerge from purely local rules, which is the same claim that
decentralized reinforcement learning makes. However, it was built for
visual plausibility, and not for coverage. Left unmodified, its
cohesion rule is harmful, due to the fact that agents which stay
together will sweep the same cells. Therefore any coverage use will
require a dispersion term. Recent
surveys~\cite{xu2022,alqudsi2024,arshid2025} treat bio-inspired
methods as the standard baseline for learned ones. Shang et
al.~\cite{shang2025} come closest to a bio-inspired coverage method.
However, their objective represents exploration efficiency, and not
coverage at a fixed budget.

\textbf{Artificial potential fields.} Khatib~\cite{khatib1986} treats
the robot as a particle under attraction to its goal and repulsion
from obstacles. It is fully local, it requires no map, and it will
cost constant time per step. Therefore it represents the natural
classical comparator for a learned local policy, and
Choset~\cite{choset2001} places it within the wider coverage
literature. Its known weakness is the local minimum in concave
geometry. The escape remedy that is utilized will strongly determine
measured performance, and papers will routinely under-specify it. Thus
an APF baseline with a stated escape rule and a published sweep
represents a harder comparator than the literature usually provides.
Hickling et al.~\cite{hickling2025} show how the field now treats APF,
due to the fact that they utilize an APF-derived reward to train a TD3
policy, rather than utilizing APF as a controller. In addition,
Du~\cite{du2023} provides a classical planning comparator based upon
an enhanced A$^*$.

\textbf{MARL for coverage, and the credit-assignment problem.}
Schulman et al.~\cite{schulman2017} introduced PPO, which now
represents the default policy-gradient method. Fusco et
al.~\cite{fusco2025} state the central difficulty for the multi-agent
case, i.e., that the joint action space will grow as $|A|^N$, which
they call "impractical" for any significant number of agents. In
addition, they state that full decentralization will scale, but that
it will produce suboptimal actions and will converge poorly under high
partial observability. They report training time rising from 0.065\,s
at $16\times10^3$ cells to 3{,}354\,s at $64\times10^3$.

There are two standard responses, and this project utilized both of
them in turn. 1. The first will sidestep the joint action space with
hierarchy or assignment, which is what most empirical work does. Each
such choice will buy tractability by reintroducing something central.
Yu et al.~\cite{yu2025} utilize event-triggered hierarchical RL, where
a high-level module will assign regions and low-level policies will
execute inside them. The assignment represents a central decision.
Therefore the result speaks to hierarchical coordination, rather than
to what independent agents will achieve. Leong et al.~\cite{leong2024}
learn selection over a discretized set of candidate viewpoints, which
bounds the action space. However, it presupposes a generator with
global knowledge of the map. Xie et al.~\cite{xie2025} decompose and
allocate cells centrally, which is efficient but is not decentralized
in the sense at issue here. None of the three establishes whether
agents acting on local input alone can cover an unknown area, due to
the fact that in each case the hardest coordination decision is taken
somewhere the agents are not. 2. The second response is
\emph{parameter sharing}, in which one policy will be applied
independently to every agent, so the action space stays
$\text{Discrete}(|A|)$ while all agents learn. Only the second will
leave the question intact, and therefore it represents the route that
was taken here.

Parameter sharing alone will not be sufficient, due to the fact that a
shared policy trained on a shared team reward will suffer the
credit-assignment problem. An agent's reward is dominated by the
behavior of teammates that it does not control. Therefore the
correlation between its own action and its own return is small, and it
will learn very little. The pathology is well documented under the
name of the "lazy agent", and there are two standard responses: 1.
difference rewards and counterfactual baselines, as in
COMA~\cite{foerster2018}; and 2. value decomposition, as in
VDN~\cite{sunehag2018} and QMIX~\cite{rashid2018}. This literature is
directly relevant here, due to the fact that
Section~\ref{sec:hybrid_ablation} shows that the earlier hybrid failed
for exactly this reason. In addition, Section~\ref{sec:reward} adopts
the simplest effective remedy (i.e., each agent will be rewarded for
the cells that it is first to visit).

\textbf{Evaluation practice, and why numerical comparison is not
available.} Replication counts across these papers are typically under
ten, where they are stated at all. None will report a hypothesis test
between its method and its baselines, none will report an effect size,
and train/test separation over layouts is rare. Therefore a
like-for-like coverage comparison is not possible. Fusco et al.\ will
report no coverage percentage at all, due to the fact that their
environment terminates on full exploration. Thus coverage represents
100\% whenever an episode completes, and their metrics are energy,
mission time, and reward. The problem recurs for three reasons. 1.
Papers will measure different quantities. 2. They will utilize
different termination rules, so fixed-budget and time-to-completion
framings represent inverses of one another. 3. They differ in scale by
more than an order of magnitude, and coverage is strongly sensitive to
the agent-to-area ratio. Therefore Table~\ref{tab:related} compares
experimental design and evidentiary standard, rather than headline
numbers.

\begin{table}[htb]
\centering
\caption{Experimental Design Compared Against Related Work.
n/s = not stated.}
\label{tab:related}
\scriptsize
\setlength{\tabcolsep}{3pt}
\begin{tabular}{llccc}
\toprule
\textbf{Study} & \textbf{Coverage metric} & \textbf{Dens.} &
\textbf{Seeds} & \textbf{Sig.\,/\,ES} \\
\midrule
Reynolds~\cite{reynolds1987} & none (animation) & no & --- & no \\
Fusco~\cite{fusco2025}       & none; energy, time & no & 1 & no \\
Yu~\cite{yu2025}             & 3D region coverage & no & n/s & no \\
Leong~\cite{leong2024}       & viewpoint coverage & no & n/s & no \\
Xie~\cite{xie2025}           & reduction in grids & no & n/s & no \\
Shang~\cite{shang2025}       & exploration effic. & no & n/s & no \\
Hickling~\cite{hickling2025} & distance, makespan & no & n/s & no \\
Du~\cite{du2023}             & path length & no & n/s & no \\
\midrule
\textbf{This work} & \textbf{\% free cells and absolute cells} &
\textbf{4} & \textbf{50} & \textbf{yes} \\
\bottomrule
\end{tabular}
\end{table}

%% ===================== METHODOLOGY =============================
\section{METHODOLOGY}\label{sec:methodology}

\subsection{System Architecture}\label{sec:arch}

Fig.~\ref{fig:sysarch} shows how the pieces will fit together. A
single world holds the obstacle map, the shared visited map, and the
six UAV positions. It will advance only when an action has been
supplied for every agent. Therefore all six will move simultaneously
within a timestep. The same world class serves all four controllers,
and the same seed will produce the same obstacle layout for each of
them. That is what makes the comparison "matched" rather than merely
parallel.

\begin{figure*}[t]
    \centering
    \includegraphics[width=0.84\textwidth]{figures/results/system_architecture.png}
    \caption{System architecture. One world serves all four controllers
    on identical seeded layouts, advancing only once all six actions
    are supplied.}
    \label{fig:sysarch}
\end{figure*}

\subsection{Environment Definition}\label{sec:env}

The arena represents a $50\times50$ discrete grid, which was built in
Python with Mesa~\cite{mesa}. At 3~m per cell the grid will cover
$150\,\text{m}\times150\,\text{m}$ (i.e., approximately a residential
block). Six UAVs start spread around a shared HOME point at grid
center $(25,25)$, eight cells out on the four cardinals and
$(\pm5,\pm5)$ on two diagonals, which reflects deployment from one
launch site. In addition, a five-cell safe zone around HOME will be
kept clear of obstacles.

Early versions scattered single-cell obstacles at random, which
produced environments that were nothing like a real search area.
Therefore the generator was rewritten. It will now place rectangular
building blocks three to eight cells on a side, with walls and L- and
U-shaped structures, at four densities (i.e., 0\%, 10\%, 20\%, and
30\% of grid area). Fig.~\ref{fig:arena} shows one such layout at
20\%, with the swarm part-way through an episode.

\begin{figure}[htb]
    \centering
    \includegraphics[width=\linewidth]{figures/methodology/scenario_arena.png}
    \caption{The task. Six UAVs sweep a
    $150\,\text{m}\times150\,\text{m}$ block at 20\% density, at step
    110 of 700 with 30\% covered. Layout and paths are both real: the
    deployed policy on a layout the experiments ran. Paths zigzag
    because actions are sampled, not greedy
    (Section~\ref{sec:eval}). Each agent sees only a window like the
    one drawn, and nothing joins one agent to another.}
    \label{fig:arena}
\end{figure}

Each episode will run 700 timesteps for every controller. Read against 3~m cells at one cell per step, that will scale to approximately 3~m/s and a 12-minute mission, which sits inside the endurance of a typical multirotor. However, that scaling represents an order-of-magnitude check and no more. The model is purely kinematic, with no inertia, turn radius, acceleration limit, or wind. In addition, diagonal moves will cover $\sqrt{2}$ cells per timestep, so implied speed will vary with heading. The figures show that the episode length is not absurd. However, they do not support a flight-time claim.
Coverage represents the fraction of non-obstacle cells visited at least once,
\begin{equation}\label{eq:coverage}
    C = \frac{|\mathcal{V}|}{|\mathcal{F}|} \times 100\%
\end{equation}
with $\mathcal{V}$ the visited free cells and $\mathcal{F}$ all free
cells. Section~\ref{sec:absolute} will report $|\mathcal{V}|$ directly
alongside $C$, due to the fact that the denominator changes with
density.

\subsection{Justification of the 90\% Target}\label{sec:threshold}

The criterion is based upon search and rescue planning, where the
governing quantity represents probability of detection (i.e., a
function of a coverage factor $C_{f} = n_{s}dW/A$ for $n_s$ searchers
traveling distance $d$ with sweep width $W$ over area $A$).
Perkins~\cite{perkins2013} states that a probability of detection of
90\% will require $C_f = 1.2$, and he records that standard tactics
will reach 63\%, 72\%, and 87\%. None of them reaches 90\%. Therefore
the value represents a recognized and demanding planning target, and
the same framework underpins the US National Search and Rescue
Supplement~\cite{nss2018}. However, the construct differs from the one
that is measured here, due to the fact that Eq.~\ref{eq:coverage} will
count cell entry with a sensor that never fails. Thus it represents an
upper bound on field performance. The criterion was fixed before any
experiment was run, and it was not adjusted afterwards.

\subsection{Observation Space}\label{sec:obs}

Each agent will sense a $7\times7$ window centered on itself (i.e., a
detection radius of three cells, or 9~m), which is realistic for a
depth sensor such as the Intel RealSense L515. The window will encode
obstacles as $-1.0$, unvisited free cells as $0.0$, visited cells as
$1.0$, and other agents as $2.0$. Therefore it provides 49 values,
plus two for normalized position (Fig.~\ref{fig:obs_window}). Writing
$p_i=(r_i,c_i)$ for agent $i$'s cell and $\delta\in\{-3..3\}^2$ for an
offset within the window,
\begin{equation}\label{eq:obs}
o_i[\delta]=
\begin{cases}
-1 & p_i{+}\delta \text{ obstacle or off-grid}\\
\;\,2 & \text{occupied by another agent}\\
\;\,1 & \text{visited by any agent}\\
\;\,0 & \text{free and unvisited}
\end{cases}
\end{equation}
with $o_i \in \mathbb{R}^{60}$ after appending $(r_i/G, c_i/G)$ and
the nine memory features of Section~\ref{sec:memory}. Cells outside
the grid will also be $-1$. Therefore an agent cannot distinguish the
world boundary from an obstacle.

The visited channel represents the swarm's shared record, and not the
agent's private one, due to the fact that a cell will read $1.0$ once
any agent has occupied it. An agent can tell from its window whether
the ground ahead was already swept. However, it cannot tell by whom or
when. This represents indirect coordination through a shared "trace",
and on a real platform it would be provided by a shared occupancy
grid, or by visible evidence of a prior pass. No agent will send or
receive a message. Therefore the property under test is decentralized
control without message passing. Complete informational isolation
represents a stronger condition, and this design does not meet it. In
addition, the channel is identical for Boids and APF, so it will
advantage no one.

\begin{figure}[htb]
    \centering
    \includegraphics[width=0.86\linewidth]{figures/methodology/fig_uav_observation_window.png}
    \caption{The $7\times7$ observation window. Values are the same
    for every controller; only the learned policy additionally
    receives the nine memory features of Section~\ref{sec:memory}.}
    \label{fig:obs_window}
\end{figure}

\subsubsection{Onboard memory features}\label{sec:memory}

The sensor window will provide no gradient towards unexplored ground
once the immediate neighborhood is swept, and a policy without one
will perform no better than a random walk. Therefore each agent will
maintain its own belief map of cells that it has sensed, and it will
derive nine features from it. These are as follows: 1. four quadrant
scores for remaining unexplored area; 2. a normalized vector and
distance to the nearest cell that it believes to be unexplored; 3. the
fraction of the arena that it has sensed; and 4. the fraction of the
episode remaining. Two of the nine (i.e., the sensed fraction and the
time remaining) are internal to the agent. The other seven will be
derived from a target set that combines what this agent has sensed
with the same shared visited record that the window already carries.
Therefore they will extend its range, but they will introduce no
mechanism that the window does not already utilize. They correspond to
onboard map memory, which any real UAV will carry. The observation
represents 60 values in total. Section~\ref{sec:ablation} will measure
what these features are worth, due to the fact that they represent
information which the classical controllers do not have. Thus the
comparison against those controllers is not one of equal information.

\subsection{Action Space}\label{sec:action}

Each UAV will move one cell per timestep in one of eight directions
(i.e., a first-order kinematic model with no inertia), which provides
$\text{Discrete}(8)$ per agent. Obstacle avoidance will happen at the
planning layer, due to the fact that cells containing obstacles are
removed from the legal move set and each controller will then choose
among the remainder. In addition, inertia, wind, and battery drain
beyond the mission budget are not modeled.

\subsection{Reward Function and Gain Selection}\label{sec:reward}

Agent $i$ receives
\begin{equation}\label{eq:reward}
    r_{i,t} = w_{cov}\,\Delta v_{i,t} - w_{time}
              + \mathbb{1}[\text{done}]\;w_{team}\,C_T/100
\end{equation}
where $\Delta v_{i,t}$ will count the cells \emph{agent $i$ was first
to visit} at step $t$, and $C_T$ represents final swarm coverage.
Rewarding own-first-visits, rather than the global new-cell count, is
what makes the signal a function of the agent's own action. Therefore
Section~\ref{sec:hybrid_ablation} will show what happens without it.

The three gains were set by scale reasoning rather than by search, and
the reasoning is worth stating, due to the fact that the relative
magnitudes will matter more than the absolute values
(Fig.~\ref{fig:reward}).

\begin{figure}[htb]
    \centering
    \includegraphics[width=\linewidth]{figures/results/reward_structure.png}
    \caption{What each reward term contributes to one agent over one
    episode, on a symmetric log scale. The coverage term dominates by
    more than an order of magnitude; the other two shape behaviour
    without competing with it.}
    \label{fig:reward}
\end{figure}

$w_{cov}=3.0$ and $w_{time}=0.05$ will set the exchange rate between
finding new ground and spending time (i.e., one new cell is worth 60
steps of the time penalty). This must be large, due to the fact that
an agent which trades a new cell for a shorter path will defeat the
objective. However, it must also be finite, since otherwise a policy
that stalls indefinitely near unexplored ground will incur no cost.
Sixty steps represents approximately a tenth of the episode, which
makes loitering expensive without allowing the penalty to compete with
the coverage signal. Over a full episode the coverage term will
accumulate to approximately $750$, against $-35$ from the time
penalty. Therefore the ratio is deliberately "lopsided".

$w_{team}=2.0$ will apply once, at termination, scaled by final swarm
coverage. Its purpose is to retain a shared interest, so that agents
will not compete for the same cells. Its magnitude is set so that it
cannot dominate, i.e., at most $2.0$ against an accumulated $\approx
750$, which represents about $0.3\%$ of return. Therefore it
represents a "tie-breaker" between otherwise equivalent policies, and
not a driver. Setting it much larger will reintroduce the
credit-assignment problem that the per-agent term exists to solve, due
to the fact that a large shared terminal payment is again mostly
determined by teammates.

\subsection{PPO Architecture}\label{sec:ppo}

All six UAVs will be driven by one shared policy, which acts on each
agent's own local observation (Fig.~\ref{fig:marlarch}). No agent is
hand-coded, and no agent will send a message.

This represents centralized training with decentralized execution.
Training will pool all six agents' experience into one gradient
update, which could not be done onboard. However, execution is local.
Each agent will run its own copy of the network on its own $7\times7$
observation, nothing is transmitted, and no controller is in the loop.
A deployed swarm will require the shared network to be installed
before flight, and it will be independent after that. In addition,
parameter sharing will keep the action space at $\text{Discrete}(8)$
rather than $8^6 = 262{,}144$, which represents the "scaling wall"
that Fusco et al.~\cite{fusco2025} describe. A naive joint-action
formulation hit it directly during this project, due to the fact that
the policy collapsed onto a single repeated action and never
recovered.

\begin{figure}[htb]
    \centering
    \includegraphics[width=0.90\linewidth]{figures/results/marl_architecture.png}
    \caption{Parameter-shared multi-agent PPO. One network produces an
    action for each agent from that agent's own observation, and each
    agent's experience is a separate sample for the same update.}
    \label{fig:marlarch}
\end{figure}

The network represents a $[256,256]$ MLP with separate policy and
value heads (i.e., the Stable-Baselines3~\cite{sb3} \texttt{MlpPolicy}
default). Six agents in each of eight parallel worlds will present 48
agent-slots to the learner. Therefore one policy update will draw upon
all six agents' experience simultaneously.

\subsection{Training Procedure}\label{sec:training}

Algorithm~\ref{alg:train} states the loop. The centralized step
represents line~\ref{ln:update}, which will pool all $NW$ agent-slots
into one gradient update. Nothing else will cross between agents. In
addition, at execution only lines~\ref{ln:obs}--\ref{ln:act} will run,
on each agent's own copy of $\pi_\theta$.

\begin{algorithm}[htb]
\caption{Parameter-shared MARL training (CTDE)}
\label{alg:train}
\begin{algbody}
\alplain{\alkw{Input:} density $d$, budget $T$, seed pool
  $\mathcal{S}_A$, $W$ worlds, $N$ agents}
\alplain{\alkw{Output:} one shared policy $\pi_\theta$}
\vspace{1pt}\hrule\vspace{2pt}
\albox{initialise $\pi_\theta$, a $[256,256]$ MLP shared by all agents}
\albox{$w_1..w_W \leftarrow$ worlds at density $d$, seeds from
  $\mathcal{S}_A$}
\albox{$t \leftarrow 0$}
\albox{\alkw{while} $t < T$ \alkw{do}}
\albox{\alind $\mathcal{B} \leftarrow \emptyset$}
\albox{\alind \alkw{for} $n_{steps}$ \alkw{do}}
\albox{\alind\alind \alkw{foreach} world $w$, agent $i \in 1..N$
  \alkw{do}}
\albox{\label{ln:obs}\alind\alind\alind $o_i \leftarrow
  \textsc{Observe}(w,i)$ \alcom{Eq.\,\ref{eq:obs}}}
\albox{\label{ln:act}\alind\alind\alind $a_i \sim
  \pi_\theta(\cdot \mid o_i)$ \alcom{sampled, not greedy}}
\albox{\alind\alind $\mathbf{r},\,done \leftarrow
  \textsc{Step}(w,\mathbf{a})$ \alcom{all $N$ move together}}
\albox{\alind\alind $\mathcal{B} \leftarrow \mathcal{B} \cup
  \{(o_i,a_i,r_i)\}_{i=1}^{N}$ \alcom{$r_i$ own-first, Eq.\,\ref{eq:reward}}}
\albox{\alind\alind \alkw{if} $done$ \alkw{then} reset $w$ with a fresh
  seed}
\albox{\label{ln:update}\alind $\theta \leftarrow
  \textsc{PPO-Update}(\theta, \mathcal{B})$
  \alcom{pools all $NW$ slots}}
\albox{\alind $t \leftarrow t + N W n_{steps}$}
\albox{\alkw{return} $\pi_\theta$}
\end{algbody}
\end{algorithm}

Hyperparameters:
$n_{steps}=64$ per slot (i.e., 3{,}072 samples per update), batch 512,
four epochs, $\gamma=0.99$, $\lambda=0.95$, clip 0.2, entropy 0.005,
learning rate $5\times10^{-4}$, one million agent-steps per density,
$W=8$, and $N=6$. Training layouts will be drawn from Seed Set A only.


Rollout length represented the decisive hyperparameter. With 48
agent-slots, $n_{steps}=256$ will yield 12{,}288 samples per update.
However, it provides only approximately forty gradient updates per
500k steps, which is far too few for PPO to converge. Therefore
coverage will stall near the random-walk level of 50\%, and it will
not improve with further training. Shortening the rollout to 64
provides approximately 160 updates over the same sample budget, and
that is what makes the task learnable. The effective number of updates
will fall as the slot count rises. Thus anyone scaling parameter
sharing to many agents will have to compensate for it.

\textbf{Training cost.} On four CPU threads of an Apple M1 the
implementation will sustain approximately 11{,}400 agent-steps per
second. Therefore one million will take 1.5 minutes, and the
four-density set approximately six, with no GPU. Fusco et
al.~\cite{fusco2025} report 3{,}354\,s for one run at $64\times10^3$
cells. The environments differ too much for a direct comparison.
However, parameter sharing will evidently keep a decentralized policy
within minutes of laptop compute at this scale. In addition, the
classical controllers will require no training at all, and that
asymmetry belongs in any deployment decision.

\subsection{Evaluation Procedure}\label{sec:eval}

Sampling at line~\ref{ln:act} represents a property of the deployed
controller, and not exploration noise that was left switched on. Six
agents will share one network. Therefore a deterministic argmax will
make every agent in an identical local situation choose the identical
action, which collapses the swarm into one repeated move and costs
almost ninety points of coverage when it is measured directly.
Sampling will break the symmetry between otherwise indistinguishable
agents, and it represents the mechanism by which one shared network
will produce six behaviors. A fielded swarm would run the same
stochastic policy. In addition, sampling is seeded from the episode
seed, so a seed will reproduce the identical episode.

\subsection{Classical-Controller Implementation}\label{sec:classical}

\textbf{Boids}~\cite{reynolds1987} will score neighboring cells as $s
= w_{cov}c_{cov} + w_{sep}c_{sep} + \mathcal{U}(-0.1,0.1)$, where
$c_{cov}$ equals 2 for an unvisited cell and $-1$ for a visited one,
and $c_{sep}$ represents the normalized distance to the nearest agent.
Agents will move sequentially within a timestep. Therefore each will
see where the others have gone, which produces spatial separation
without communication.

\textbf{APF}~\cite{khatib1986} scores a candidate cell as
\begin{equation}\label{eq:apf_score}
    s = c_{cov}
      - w_{rep} \!\!\sum_{o \in \mathcal{O}_{n}} \frac{k_{rep}}{d_o^2}
      + c_{sep} + \mathcal{U}(-0.1, 0.1)
\end{equation}
over obstacles within radius $r_{rep}$, with $c_{cov}$ of 5.0 for a
cell unvisited by any agent, 1.0 for one absent from the agent's own
memory, and $-2.0$ otherwise, and $w_{rep}=0.3$.
Section~\ref{sec:sweep} will sweep the repulsion term of
Eq.~\ref{eq:apf_score}. An escape rule will fire a random move to an
unvisited neighbor after three stalled steps, and that is what
prevents the local-minimum collapse for which the method is known.
Rather than assert that $k_{rep}=1.5$ and $r_{rep}=2$ are well chosen,
Section~\ref{sec:sweep} will report the coverage surface over both.

\subsection{Experimental Design and Statistical Methodology}
\label{sec:stats}

The experiment crosses four controllers with four densities and two
seed sets (i.e., 32 conditions, 50 seeds each, and 1{,}600 runs). Seed
Set~A (0--49) was utilized for development and training, and Seed
Set~B (50--99) was reserved and utilized only at the end. Every
controller will see the identical 50 layouts within a cell, as
Algorithm~\ref{alg:eval} states. Each will pass its seed to the same
obstacle generator, and the layouts are regenerated and compared cell
by cell before any test runs. Therefore this is what licenses the
paired design. The implementation represents Python~3.13 with
Mesa~\cite{mesa}, Stable-Baselines3~\cite{sb3}, and Gymnasium. In
addition, the bootstrap is seeded at 20260729, and code, results, and
checkpoints will accompany this report.

\begin{algorithm}[htb]
\caption{Matched-layout evaluation}
\label{alg:eval}
\begin{algbody}
\alplain{\alkw{Input:} controllers $\mathcal{C}$, densities
  $\mathcal{D}$, seed sets $\mathcal{S}_A,\mathcal{S}_B$}
\vspace{1pt}\hrule\vspace{2pt}
\albox{\alkw{foreach} $d \in \mathcal{D}$, $\mathcal{S} \in
  \{\mathcal{S}_A,\mathcal{S}_B\}$, seed $s \in \mathcal{S}$ \alkw{do}}
\albox{\alind $M_s \leftarrow \textsc{GenObstacles}(d,s)$;
  \alkw{assert} identical for every $c \in \mathcal{C}$}
\albox{\alind $C[c,s] \leftarrow \textsc{Run}(c,M_s,700)$
  \alkw{for each} $c \in \mathcal{C}$}
\albox{\alkw{foreach} pair $(c_1,c_2)$, each $(d,\mathcal{S})$
  \alkw{do}}
\albox{\alind Wilcoxon on $\{C[c_1,s]-C[c_2,s]\}_s$; $r_{rb}$,
  Cliff's $\delta$, bootstrap CI}
\albox{Holm--Bonferroni within each comparison family}
\end{algbody}
\end{algorithm}

There are three measures in Table~\ref{tab:extended} that each admit
more than one reasonable definition, and every controller will be
measured on the same one. With $u_i$ the distinct cells agent $i$
occupied,
\begin{equation}\label{eq:secondary}
L=\ell\!\!\sum_{i=1}^{N}\sum_{t}\lVert p_{i,t}-p_{i,t-1}\rVert_2,
\quad
\mathcal{G}=\frac{2\sum_{k} k\,u_{(k)}}{N\sum_{k}u_{k}}-\frac{N{+}1}{N}
\end{equation}
with $\ell=3$~m per cell, $u_{(k)}$ the $k$th smallest $u_i$, and
$\mathcal{G}=0$ perfectly even. A "near-miss" represents a pair within
one cell by Euclidean distance, which will be counted before each
move.

Twenty-seven of the 32 condition-level coverage distributions will
depart from normality under Shapiro-Wilk at $\alpha = 0.05$. Therefore
rank-based tests are primary. The design is matched, and this was
verified rather than assumed, due to the fact that for each seed the
obstacle map was regenerated independently through the Mesa generator
and through both learned environments, and then compared cell by cell.
The maps are identical for every sampled seed at every density.
Therefore observation $i$ in each sample will refer to the same
environment, and the paired Wilcoxon signed-rank test is correct.
Mann-Whitney $U$~\cite{mann1947} and Welch's $t$ will be reported as
robustness checks. Holm-Bonferroni correction~\cite{holm1979} will be
applied within each family. Effect sizes represent matched-pairs
rank-biserial correlation $r_{rb}$ and Cliff's
$\delta$~\cite{cliff1993}, and intervals represent bootstrap
percentile intervals from 10{,}000 resamples~\cite{efron1979}. In
addition, all error bars in this report are 95\% confidence intervals
on the mean, rather than standard deviations.

%% ===================== RESULTS =================================
\section{RESULTS}\label{sec:results}

Table~\ref{tab:main_results} will report final coverage for every
controller at every obstacle density on both seed sets, with the
number of runs reaching the 90\% target given in parentheses.

\begin{table}[htb]
\centering
\caption{Final Coverage (\%), Mean $\pm$ SD over 50 Seeds, with Runs
Reaching the 90\% Target in Parentheses}
\label{tab:main_results}
\scriptsize
\setlength{\tabcolsep}{0.8pt}
\resizebox{\columnwidth}{!}{%
\begin{tabular}{@{}l@{\hspace{2pt}}lcccc@{}}
\toprule
\textbf{Set} & \textbf{Controller} & \textbf{0\%} & \textbf{10\%} &
\textbf{20\%} & \textbf{30\%} \\
\midrule
\multirow{4}{*}{A}
 & Boids  & $87.9\pm4.5$\,(17) & $86.7\pm4.8$\,(12) & $85.7\pm5.6$\,(12) & $84.6\pm6.6$\,(9) \\
 & APF    & $91.4\pm3.6$\,(36) & $92.5\pm3.9$\,(39) & $94.2\pm3.5$\,(42) & $\mathbf{94.7\pm2.7}$\,(47) \\
 & Hybrid & $85.8\pm5.7$\,(13) & $85.0\pm5.5$\,(9)  & $84.8\pm7.0$\,(13) & $83.8\pm6.1$\,(10) \\
 & MARL   & $\mathbf{96.1\pm1.5}$\,(50) & $\mathbf{93.7\pm2.5}$\,(45) & $93.5\pm3.7$\,(41) & $90.8\pm5.7$\,(35) \\
\midrule
\multirow{4}{*}{B}
 & Boids      & $86.6\pm4.6$\,(12) & $86.8\pm6.4$\,(18) & $87.5\pm4.9$\,(17) & $83.6\pm6.7$\,(7) \\
 & APF        & $90.5\pm4.1$\,(27) & $93.0\pm3.8$\,(41) & $\mathbf{94.5\pm2.8}$\,(47) & $\mathbf{94.0\pm4.4}$\,(42) \\
 & Hybrid     & $85.6\pm5.2$\,(11) & $85.9\pm4.7$\,(8)  & $84.4\pm6.3$\,(9)  & $82.2\pm6.8$\,(5) \\
 & MARL   & $\mathbf{95.5\pm1.8}$\,(49) & $\mathbf{93.8\pm2.7}$\,(46) & $92.8\pm3.8$\,(36) & $93.1\pm3.8$\,(41) \\
\bottomrule
\end{tabular}}
\end{table}

\begin{figure*}[t]
    \centering
    \includegraphics[width=0.84\textwidth]{figures/results/four_way_coverage_ci.png}
    \caption{Coverage against obstacle density for all four
    controllers, on the development and held-out seed sets. Error bars
    are 95\% bootstrap confidence intervals on the mean. Only MARL and
    APF clear the target, crossing between 10\% and 20\% on both sets.}
    \label{fig:fourway}
\end{figure*}

\subsection{Meeting the 90\% Criterion}

Both the MARL policy and APF will exceed 90\% by mean in all eight of
their cells. One-sided Wilcoxon tests against 90\%, Holm-corrected
over all 32 cells, will confirm that the excess is significant in
seven cells each. The exceptions represent the cells where the mean
sits closest to the line, i.e., MARL at 30\% density on the
development set (90.84\%, Holm $p = 0.499$) and APF at 0\% density on
the held-out set (90.46\%, Holm $p = 1.000$). Boids will fall
significantly short in seven cells, and it is indistinguishable from
90\% in one. In addition, the hybrid will fall short in all eight.
Therefore only the two leading controllers will meet the criterion at
all, and neither dominates the other on it.

\subsection{Pairwise Comparison}\label{sec:pairwise}

\begin{table}[htb]
\centering
\caption{MARL Against Each Baseline, Paired Wilcoxon Signed-Rank,
Holm-Corrected. Difference in percentage points.}
\label{tab:stats}
\scriptsize
\setlength{\tabcolsep}{1.6pt}
\begin{tabular}{lllrcrr}
\toprule
\textbf{Set} & \textbf{Dens.} & \textbf{Comparison} & \textbf{Diff.} &
\textbf{95\% CI} & \textbf{Holm $p$} & \textbf{$r_{rb}$} \\
\midrule
\multirow{6}{*}{A}
& 0\%  & MARL vs APF    & $+4.71$  & [3.67, 5.81]        & $<$0.0001 & $+0.94$ \\
& 10\% & MARL vs APF    & $+1.21$  & [$-0.09$, 2.53]     & 0.5523    & $+0.24$ \\
& 20\% & MARL vs APF    & $-0.65$  & [$-2.02$, 0.72]     & 0.6892    & $-0.20$ \\
& 30\% & MARL vs APF    & $-3.82$  & [$-5.59$, $-2.15$]  & 0.0004    & $-0.64$ \\
& 0\%  & MARL vs Boids  & $+8.19$  & [7.00, 9.40]        & $<$0.0001 & $+1.00$ \\
& 0\%  & MARL vs Hybrid & $+10.28$ & [8.72, 11.95]       & $<$0.0001 & $+1.00$ \\
\midrule
\multirow{6}{*}{B}
& 0\%  & MARL vs APF    & $+5.02$  & [3.72, 6.31]        & $<$0.0001 & $+0.87$ \\
& 10\% & MARL vs APF    & $+0.82$  & [$-0.59$, 2.30]     & 0.9649    & $+0.11$ \\
& 20\% & MARL vs APF    & $-1.73$  & [$-2.99$, $-0.51$]  & 0.2702    & $-0.33$ \\
& 30\% & MARL vs APF    & $-0.90$  & [$-2.45$, 0.68]     & 0.7116    & $-0.24$ \\
& 20\% & MARL vs Boids  & $+5.33$  & [3.69, 7.06]        & $<$0.0001 & $+0.77$ \\
& 20\% & MARL vs Hybrid & $+8.41$  & [6.41, 10.49]       & $<$0.0001 & $+0.90$ \\
\bottomrule
\end{tabular}

\vspace{0.3em}
{\footnotesize The full 48-comparison table, including every
MARL-vs-Boids and MARL-vs-Hybrid cell and the Mann-Whitney and Welch
robustness checks, is in the supplementary output.}
\end{table}

There are three findings that follow. 1. MARL will beat APF in open
terrain by 4.71 and 5.02 points, and both are highly significant with
near-maximal effect sizes. 2. The two cannot be separated at
intermediate density, due to the fact that at 10\% and 20\% no
comparison will survive Holm correction on either seed set. 3. APF
will keep an edge at the hardest density by 3.82 points with a medium
effect. However, this occurs only on the development set, since the
same comparison on the held-out set represents 0.90 points and is not
significant. Therefore the advantage is not robust across layout
samples. In addition, against the two weaker baselines MARL will win
everywhere by 5.3 to 10.9 points, with $r_{rb}$ from 0.75 to 1.00, and
all sixteen will survive correction.

\subsection{Generalisation}\label{sec:generalisation}

Seed Set~B consists of layouts that are distinct from Set~A, which was
verified by regenerating and comparing maps. Changes from A to B for
the MARL policy represent $+0.60$, $-0.11$, $+0.76$, and $-2.31$
percentage points at 0--30\% density. No controller will change by
more than 2.31 points, and no difference will survive Holm correction.
The policy was trained only on Set~A layouts. Therefore this
represents a genuine held-out test, and it shows no overfitting to the
training maps. In addition, at 30\% density it will perform
\emph{better} on the unseen set. At 0\% density there are no
obstacles, so the two sets represent the same environment, and any
difference there will reflect only the stochastic policy.

\subsection{Absolute Cells Against Percentage}\label{sec:absolute}


Absolute-cell accounting will resolve an ambiguity that the earlier
iteration of this project could not settle. APF's percentage coverage
will rise by 3.3 points from 0\% to 30\% obstacles, which invites the
reading that structured obstacles help it search. However, the
absolute counts will contradict that, and they were recovered over the
same 50 seeds by multiplying each run's coverage by the free area of
its layout. APF will cover 640 fewer cells at 30\% than at 0\% (i.e.,
a fall of 28\%), Boids 727 fewer (i.e., 33\%), and MARL 824 fewer
(i.e., 34\%). Therefore every controller will cover less ground as
density rises, while the free area itself shrinks from 2{,}500 to
1{,}737 cells (i.e., a fall of 31\%). APF's percentage will rise only
due to the fact that its absolute coverage falls more slowly than the
denominator. Thus no controller will become a better searcher in
cluttered terrain, and APF will simply degrade most gracefully. In
addition, percentage coverage alone represents a misleading metric
across varying free area, and reporting both represents the correct
practice.

Fig.~\ref{fig:traj_30} will show what that difference looks like on
the ground, i.e., three controllers sweeping one identical 20\%
layout, where APF's corridor-following and MARL's outward partition
are visible as search patterns rather than as summary statistics.

\begin{figure*}[t]
    \centering
    \includegraphics[width=0.84\textwidth]{figures/results/trajectories_four_way.png}
    \caption{Swarm trajectories on one identical layout at 20\%
    obstacle density. Boids doubles back in open areas; APF follows
    the corridors between buildings; the MARL policy sweeps
    outward and partitions the arena without communicating.
    $\star$ marks HOME. Equivalents at every density are in the
    supplementary material.}
    \label{fig:traj_30}
\end{figure*}

%% ===================== EXTENDED ANALYSIS =======================
\section{EXTENDED ANALYSIS}\label{sec:extended}

Coverage alone represents a poor guide to field viability. Therefore,
following supervisory feedback, further analyses were run at the same
700-step budget on 50 matched seeds. The exceptions represent the
swarm-size and hybrid sweeps, which will require a run per
configuration and will utilize 15.

\begin{table}[htb]
\centering
\caption{Secondary Measures, Mean over 50 Seeds, defined at
Eq.~\ref{eq:secondary}. Near-misses are a proximity proxy, not
observed collisions. $^\dagger$Per-step cost re-measured for all four
on one machine; only the ordering is meaningful.}
\label{tab:extended}
\scriptsize
\setlength{\tabcolsep}{3pt}
\begin{tabular}{llcccc}
\toprule
\textbf{Measure} & \textbf{Controller} & \textbf{0\%} & \textbf{10\%} &
\textbf{20\%} & \textbf{30\%} \\
\midrule
\multirow{4}{*}{Path length}
 & Boids  & 14913 & 14860 & 14819 & 14747 \\
 & APF    & 14946 & 14837 & 14723 & 14582 \\
 & Hybrid & 13927 & 12744 & 12690 & 13505 \\
 & MARL   & 15240 & 14399 & 13696 & 11950 \\
\midrule
\multirow{4}{*}{Gini}
 & Boids  & 0.067 & 0.077 & 0.087 & 0.103 \\
 & APF    & 0.029 & 0.039 & 0.050 & 0.083 \\
 & Hybrid & 0.197 & 0.228 & 0.239 & 0.244 \\
 & MARL   & 0.058 & 0.051 & 0.064 & 0.103 \\
\midrule
\multirow{4}{*}{Near-misses}
 & Boids  & 1244.4 & 1088.4 & 888.9 & 701.5 \\
 & APF    & 1245.4 & 1109.4 & 938.7 & 606.6 \\
 & Hybrid & 0.8    & 0.9    & 0.7   & 0.8   \\
 & MARL   & 47.8   & 55.6   & 50.1  & 69.6  \\
\midrule
\multirow{4}{*}{ms\,/\,step$^\dagger$}
 & Boids  & 0.085 & 0.212 & 0.327 & 0.431 \\
 & APF    & 0.358 & 0.581 & 0.652 & 0.776 \\
 & Hybrid & 0.412 & 0.418 & 0.408 & 0.392 \\
 & MARL   & 0.728 & 0.889 & 0.702 & 0.683 \\
\bottomrule
\end{tabular}
\end{table}

\textbf{Load balance.} APF will spread work most evenly, MARL second,
and Boids third. However, the hybrid is far worse, and the per-agent
breakdown explains why. At 20\% obstacles its agents will contribute
0.4\%, 19.6\%, 19.7\%, 20.9\%, 19.6\%, and 19.8\% of unique coverage.
Agent~0 represents the only one under policy control, and it will
account for 0.4\% against a balanced 16.7\%. MARL will show no such
pathology, and it will reach near-APF balance with no agent hand-coded
to do so, due to the fact that each is rewarded for its own
contribution.

\textbf{Proximity events.} A "near-miss" will proxy collision risk,
and it will count no collisions, due to the fact that the grid
provides a vehicle with no body. Therefore the measure will capture
how often a real airframe would have had to give way. MARL will record
approximately twenty times fewer than either classical controller,
while covering more ground. However, the hybrid's near-zero figure
comes from elsewhere, since when it is read together with the
load-balance result, it describes its single learned agent staying out
of the way.

\textbf{Computational cost.} Boids and APF represent $O(N^2)$ per
timestep, due to the fact that each agent will check every other. A
shared policy represents $O(N)$ (i.e., one forward pass each). The
constants will reflect that, since the classical pair will rise with
density while the learned ones will stay flat. However, at six agents
MARL is still the dearest of the four, due to the fact that a forward
pass will cost more than scanning five neighbors. Therefore the
scaling favors it, and at this size it will buy coverage with compute.

\textbf{Swarm size.} All four controllers were run at 3, 6, 9, and 12
agents at 20\% density over 15 seeds (Table~\ref{tab:swarm}). MARL and
the hybrid were trained at six agents. Therefore other sizes represent
a transfer test, and the shared-policy architecture is what makes this
possible at all, due to the fact that one network will apply to any
swarm size without retraining.

\begin{table}[htb]
\centering
\caption{Coverage and Near-Misses by Swarm Size, 20\% Obstacles,
15 Seeds. MARL and the hybrid are six-agent policies tested at other
sizes. All four columns count a near-miss identically: Euclidean
radius one cell, sampled before each step.}
\label{tab:swarm}
\scriptsize
\setlength{\tabcolsep}{3pt}
\begin{tabular}{lcccc|cccc}
\toprule
& \multicolumn{4}{c|}{Coverage (\%)} & \multicolumn{4}{c}{Near-misses} \\
$N$ & Boids & APF & Hyb. & MARL & Boids & APF & Hyb. & MARL \\
\midrule
3  & 63.5 & 68.9 & 49.4 & \textbf{73.6} & 185  & 175  & 2.2  & 5    \\
6  & 85.2 & 93.5 & 82.6 & \textbf{93.6} & 827  & 910  & 20.9 & 52   \\
9  & 93.9 & \textbf{99.0} & 92.8 & 98.9 & 2312 & 2059 & 55.7 & 200  \\
12 & 97.7 & \textbf{99.7} & 97.3 & 99.2 & 3845 & 3819 & 92.9 & 300  \\
\bottomrule
\end{tabular}
\end{table}

There are two things that stand out. 1. The MARL policy will lead at
three agents by 4.7 points over APF and by 10.1 over Boids, which
represents the condition where coordination matters most, due to the
fact that there is least redundancy to hide a poor sweep. It will
converge with APF at six and nine, and it will fall slightly behind at
twelve, where the arena is saturated and any reasonable controller
will approach the ceiling. 2. Returns will fall off past nine agents
for every controller, while near-misses will rise four- to sixfold
from six to twelve (Fig.~\ref{fig:sensitivity}). Therefore six agents
sits on the steep part of the curve, where coordination differences
will be amplified. In addition, across all sizes the learned
controllers will record twelve to thirty-four times fewer near-misses
than the classical ones.

\begin{figure*}[t]
    \centering
    \includegraphics[width=0.84\textwidth]{figures/extended/uav_count_sensitivity_all3.png}
    \caption{Coverage and near-misses against swarm size at 20\%
    density, 15 seeds, the same runs as Table~\ref{tab:swarm}. The
    learned controllers are dashed: both were trained at six agents
    and tested outside that configuration at every other size. Log
    axis on the right, where they stay about an order of magnitude
    below the classical pair throughout.}
    \label{fig:sensitivity}
\end{figure*}

\textbf{Moving obstacles.} With obstacles relocating every 50 steps at
20\% density, APF will fall from 94.0\% to 85.9\% as the moving share
goes from none to half. However, Boids will drop only from 85.9\% to
82.8\%. APF will require a field that is consistent long enough to
guide an agent through, whereas Boids will recompute each step. The
learned policies were not retrained under this perturbation. Therefore
this covers the classical pair only.

\subsection{Information Asymmetry: the Memory Ablation}
\label{sec:ablation}

The learned controller will carry nine onboard-memory features
(Section~\ref{sec:memory}) that Boids and APF do not have. Rather than
declare this as a caveat, the authors measured it. An otherwise
identical policy was trained with those features zeroed, which leaves
it exactly the information that the classical controllers will receive
(i.e., a $7\times7$ sensed window and its own position).

\begin{table}[htb]
\centering
\caption{Memory-Feature Ablation, 20\% Density, 50 Matched Seeds. The
Ablated Policy Keeps the Same 60-Dimensional Vector With the Nine
Memory Features Held at Zero}
\label{tab:ablation}
\scriptsize
\setlength{\tabcolsep}{3pt}
\begin{tabular}{lccc}
\toprule
\textbf{Configuration} & \textbf{Steps} & \textbf{Coverage (\%)} &
\textbf{$\geq$90\%} \\
\midrule
MARL, memory     & 1.00M & $93.6\pm3.7$ & 41/50 \\
MARL, memory     & 1.85M & $98.8\pm1.0$ & 50/50 \\
MARL, no memory  & 1.00M & $89.6\pm4.6$ & 31/50 \\
MARL, no memory  & 1.85M & $90.0\pm4.2$ & 29/50 \\
APF (same information) & --- & $94.2\pm3.5$ & 42/50 \\
\bottomrule
\end{tabular}
\end{table}

At a matched budget of one million agent-steps the memory features are
worth $+4.00$ percentage points (i.e., paired Wilcoxon, $p<0.0001$,
$r_{rb}=0.72$), and without them the policy will reach 89.6\% against
APF's 94.2\%. Therefore on equal information and equal training the
classical controller will win, and it will still win at 1.85M ($-4.22$
pp, $p<0.0001$).

Running both arms on to 1.85M will separate two explanations that the
single-budget comparison cannot. The memory-free policy will gain
nothing from the extra 850k steps ($+0.43$ pp, $p=0.83$), due to the
fact that it has converged. However, the memory policy will gain
$+5.26$ pp over the same interval and will reach 98.8\%. Therefore the
gap widens from 4.00 to 8.83 points ($r_{rb}=1.00$). Extra training
will not close it, which rules out the "head-start" explanation. What
the features will change is what the architecture can represent from
local sensing alone, i.e., a $7\times7$ window and a position alone
will not reach the target reliably.

\subsection{Credit Assignment Measured Directly}
\label{sec:credit}

Section~\ref{sec:reward} will reward each agent for the cells that it
is first to visit. The alternative that the literature warns against
represents the naive shared reward, in which every agent will receive
the global new-cell count. Difference rewards~\cite{foerster2018} and
value decomposition~\cite{sunehag2018},~\cite{rashid2018} exist to
repair that arrangement. In addition, it can be run here as a
controlled experiment, due to the fact that the network, budget, seed,
observation, and dynamics will all stay fixed, and only the reward
scheme will move.

The shared-reward policy will reach $63.0\pm6.3\%$ against
$93.6\pm3.7\%$ for the per-agent reward, which represents a gap of
$30.60$ percentage points ($p<10^{-14}$, $r_{rb}=1.00$), and not one
of its fifty runs will reach the target. Therefore on this task credit
assignment will decide whether the controller works at all, due to the
fact that on the shared reward it will lose to every baseline. In
addition, it accounts for the hybrid in
Section~\ref{sec:hybrid_ablation} without appeal to its architecture,
since the hybrid utilized exactly this reward.

\subsection{Training-Seed Variance and Budget}\label{sec:seedvar}

Every MARL figure that is reported above comes from one training run
at PPO seed 0. Fifty evaluation seeds will control layout variance.
However, they say nothing about variance across training runs.
Therefore three further policies were trained at seeds 1, 2, and 3
under an identical budget, and they were evaluated on the same fifty
layouts (Table~\ref{tab:seedvar}).

\begin{table}[htb]
\centering
\caption{Coverage by PPO Training Seed, 20\% Density, 50 Matched
Evaluation Seeds Each. APF Scores 94.2\% on the Same Layouts.
SD$_{\text{seed}}$ is Across the Four Means, not Across Runs}
\label{tab:seedvar}
\scriptsize
\setlength{\tabcolsep}{2pt}
\begin{tabular}{lccccccc}
\toprule
\textbf{Budget} & \textbf{s0} & \textbf{s1} & \textbf{s2} &
\textbf{s3} & \textbf{Mean} & \textbf{SD$_{\text{seed}}$} &
\textbf{$\geq$90\%} \\
\midrule
1.00M & 93.55 & 92.98 & 91.08 & 85.75 & 90.84 & 3.55 & 122/200 \\
1.85M & 98.81 & 98.87 & 96.81 & 94.69 & 97.30 & 1.98 & 193/200 \\
\bottomrule
\end{tabular}
\end{table}

At 1.00M the spread is wide, i.e., 7.8 points between the best and
worst seed, and seed 0 represents the best of the four. Therefore the
headline MARL numbers will sit at the top of their distribution, and
not in the middle of it. A single fresh training run at 20\% density
will average 90.8\%, with a standard deviation of 5.0 points across
the 200 individual runs. That will put it behind APF.

The 1.85M row will identify the cause. The across-seed standard
deviation will fall from 3.55 to 1.98 points, the worst seed will rise
to 94.7\% and will overtake APF, and every seed will clear the target
on at least 44 of 50 layouts. Therefore both the shortfall and the
instability at 1.00M come from stopping on the steep part of the
learning curve. At 1.85M the same policy will beat APF by 4.61 points
at 20\% density ($p<10^{-14}$, $r_{rb}=1.00$). Thus what
Table~\ref{tab:stats} records at that density represents the training
budget talking, and not a ceiling on the method. Holding the budget
fixed across every controller and density represents the right control
for the main comparison. However, at 20\% it will understate the
learned controller.

How far that will generalize is untested. The extended-budget evidence
comes from one density. In addition, 10\% represents a different
regime, with a smaller starting margin and sparser structure for APF
to exploit. The same direction is plausible there. However,
"plausible" represents an inference, and not a measurement. Therefore
repeating the protocol at every density represents future work.

\subsection{APF Parameter Sensitivity}\label{sec:sweep}

The repulsion parameters were swept over $k_{rep} \in \{0.5, 1.0, 1.5,
2.5, 4.0\}$ and $r_{rep} \in \{1,2,3,4\}$ at every density. The
reported configuration $(1.5, 2)$ will sit on a broad "plateau", due
to the fact that coverage will vary by under two points across the
center of the grid, and it will degrade sharply only at $r_{rep}=4$,
where the repulsion radius approaches the corridor width. Therefore
APF's showing represents a property of the method here, and not of a
fortunate parameter choice, which is what makes it a fair comparator.
The full surface is in the supplementary output.

%% ===================== VALIDATION ==============================
\section{VALIDATION AND THREATS TO VALIDITY}\label{sec:validity}

There were three defects found by auditing the experiment code, and
all of them were repaired before any result above was generated. In
addition, one earlier experiment is withdrawn. A fourth item (i.e.,
the ablation in Section~\ref{sec:hybrid_ablation}) will be reported as
a finding rather than as a defect.

\subsection{The Hybrid Could Not Answer the Question}
\label{sec:hybrid_ablation}

An earlier iteration utilized one policy-controlled agent alongside
five hand-written greedy agents, which represents an architecture
common enough to be worth reporting on. It carried the shared-reward
defect that is quantified in Section~\ref{sec:credit}, and a second
one that is fatal on its own. Replacing the learned agent with an
inert, a uniformly random, a greedy, and the trained agent will
provide swarm coverage of 82.21\%, 82.37\%, 84.07\%, and 84.83\% at
20\% density over 15 seeds. The trained policy represents the best of
the four, so it had learned something. However, the entire range
spanned by completely different behaviors for that one agent
represents 1.86 percentage points, against a 9.37-point gap to APF.
Therefore no agent could have closed it. Thus the hybrid's deficit
represents a property of its five reactive agents, and not of the
learned one, and no conclusion about PPO will follow from it.

The fully decentralized architecture in Section~\ref{sec:ppo} will
remove both problems. In addition, the 10.25-point margin between MARL
and the hybrid in Table~\ref{tab:stats} will quantify what the
shortcut cost.

\subsection{Repaired Defects}

\textbf{Uncontrolled obstacle seeding.} The Mesa models originally
drew the obstacle seed from Mesa's internal generator, and not from
the seed supplied by the runner. Therefore one seed gave different
layouts on repeat (i.e., 503, 510, and 512 obstacle cells across three
runs). Both models will now take an explicit seed, verified cell by
cell, and that is what licenses the paired tests.

\textbf{Seed Set B was not held out.} The earlier PPO environment
mapped the episode seed with \texttt{s = \_ep\_seed \% 50}. Therefore
Set B resolved to Set A's fifty layouts. The modulo was removed, and
the sets are now disjoint. \textbf{Unequal episode budgets.} Boids ran
for 800 steps against 700 for the others. All four will now utilize
700.

\textbf{A retracted ablation figure.} An earlier draft reported the
memory-free policy at 95.0\% after 1.85M steps, and it concluded that
the memory features bought convergence speed alone. However, the
surviving checkpoint will re-evaluate to the 1.00M figure exactly
(i.e., 89.55\%), and no checkpoint matching 95.0\% exists. Retrained
from 1.00M to 1.85M it will score $90.0\pm4.2\%$, which represents the
value in Table~\ref{tab:ablation}. Therefore the conclusion reverses.

\subsection{Withdrawn and Outstanding}

The sensor-noise experiment is withdrawn. For Boids and APF the
perturbation corrupted the obstacle grid itself, once per episode,
which marked approximately 398 of 1{,}991 free cells unreachable and
imposed a mechanical ceiling near 80\% before the controller acted at
all. However, for the hybrid it entered only the observation of the
single learned agent, which will contribute 0.4\% of coverage.
Therefore the comparison represents two different interventions, and
it will support no conclusion. A valid version will require the same
corruption, representation, and persistence for every controller.

Redundancy is recorded for the MARL controller, but not
retrospectively for the classical ones. Therefore no cross-controller
comparison will be made.

%% ===================== DISCUSSION ==============================
\section{DISCUSSION}\label{sec:discussion}

\textbf{Answering the question.} A decentralized multi-agent PPO
policy (i.e., one in which every UAV learns from a local window,
executes on its own network, and exchanges no messages) will reach
90\% coverage at every obstacle density tested, and on layouts that it
has never seen. APF will do so too, on the same layouts and by the
same test. However, the two learned-free baselines will not. Against a
well-tuned APF baseline it will win clearly in open terrain, it will
tie at intermediate density, and at the hardest density it will trail
on the development set but not on the held-out one. Therefore the
answer to the research question represents "yes" for the criterion,
and "qualified" for the comparison, i.e., decentralized learning will
match classical reactive control on this task, and it will beat it
where the geometry offers least structure.

\textbf{Why the crossover occurs.} The two controllers will fail in
opposite regimes, and the crossover in Fig.~\ref{fig:fourway}
therefore represents the most informative result here. In an empty
grid APF will have little to work with, due to the fact that six
agents under mutual repulsion will spread evenly and will then have
nothing to distinguish one direction from another, so they will
re-sweep ground. However, the learned policy will carry a directional
signal from its own belief map, and it will keep moving towards
unexplored space, which represents exactly what an empty arena
rewards. As obstacles increase, the situation narrows. Walls and
corridors will provide APF's repulsion term with the structure that it
requires, and they will channel agents along streets and decorrelate
their paths. Meanwhile the learned policy will face a harder
navigation problem, with a 7$\times$7 window that cannot see around a
building. The 30\% result is significant on one seed set and not on
the other, which suggests that the remaining gap is small relative to
layout-to-layout variation. Therefore the classical controller will
degrade most gracefully in clutter, while the learned one will exploit
open space best.

Section~\ref{sec:seedvar} will limit how far that reading extends. The
crossover is measured at a fixed budget, and at 1.85M the learned
controller will reach 98.8\% at 20\% density and will overtake APF
there on every training seed. Therefore what the main experiment
establishes is that the two are close under equal training, and not
that clutter will impose a ceiling that the learned controller cannot
pass.

\textbf{What the absolute-cell result changes.} No controller will
improve in absolute terms as obstacles increase, which is contrary to
what percentage coverage suggests. This matters beyond the present
study, due to the fact that coverage over free cells represents the
standard metric across the literature surveyed in
Section~\ref{sec:related}. Therefore any paper that varies obstacle
density while reporting only that ratio will invite the same
misreading.

\textbf{Relation to the literature.} Fusco et al.~\cite{fusco2025}
predict that fully decentralized MARL will underperform under high
partial observability, and that the joint action space will make the
design intractable at scale. This project confirms the second half
directly, due to the fact that it hit the $8^6$ wall. However, it
qualifies the first, since with parameter sharing, a per-agent reward,
and onboard memory, a fully decentralized policy is competitive with
(and in open terrain better than) a tuned classical controller.
Therefore the result suggests that the pessimism concerns the naive
formulation, rather than decentralization itself. The
credit-assignment
literature~\cite{foerster2018,sunehag2018,rashid2018} represents the
reason, and Section~\ref{sec:credit} will put a number on it, i.e.,
the same network on the shared team reward that those papers set out
to repair will lose 30.6 points and will never reach the target. In
this project, the difference between a policy that fails and one that
beats APF was mostly a matter of what the agent was rewarded for.

\textbf{Practical implication.} For open or lightly cluttered search
areas the learned controller will cover most ground, and it will
record between nine and twenty-six times fewer proximity events than
either classical controller, with the margin narrowing as density
rises. However, for dense urban clutter a tuned potential field will
clear the same target, it will remain slightly ahead on coverage, and
it is far cheaper to deploy, due to the fact that it requires no
training and no GPU. Therefore the learned controller is worth its
cost only if it is trained past the point where its across-seed spread
closes. In addition, for any mission, adding agents will rival
changing algorithm, i.e., six to nine UAVs will gain 5.3 to 10.2
points, against an 11.0-point spread between the best and worst
controller at six.

%% ===================== CONCLUSION ==============================
\section{CONCLUSIONS AND FUTURE WORK}\label{sec:conclusion}

A parameter-shared multi-agent PPO policy driving six UAVs from local
observations, with decentralized execution and no message passing,
achieved 90.8\% to 96.1\% coverage across obstacle densities from 0\%
to 30\% on both development and held-out layouts. Its mean will exceed
the 90\% search-and-rescue planning target in every condition tested,
and it will do so significantly in seven of eight after correction.
APF will match that record, and the other two controllers will never
reach the target. It beat a tuned APF controller by 4.7 and 5.0
percentage points in open terrain, and it was statistically
indistinguishable at 10\% and 20\% density. However, at 30\% it
trailed by 3.8 points on the development set, which represents a
difference that did not replicate on the held-out set. In addition, it
beat a Boids baseline and a single-agent hybrid by 5.3 to 10.9 points
throughout. Transfer to fifty unseen layouts changed coverage by at
most 2.31 points, and it was never significant.

There are two design decisions that account for that performance, and
the ablations will size both. 1. Replacing the per-agent reward with
the shared team reward, and changing nothing else, will cost 30.6
points, and it will leave none of fifty runs at the target. 2.
Removing the memory features will cost a further 4.0 points at one
million steps, and the memory-free policy is still at 90.0\% after
1.85M, having gained nothing from the extra 850k. Therefore to that
budget those features will raise the ceiling, and not the speed of
reaching it. Neither decision represents a refinement, and they will
not fail alike. On the shared reward the policy will collapse to
63.0\% and will lose to every baseline here, whereas without memory it
will hold 89.6\%, ahead of Boids and the hybrid but short of both the
target and APF. In addition, that second decision will qualify the
headline, due to the fact that the memory features represent
information which Boids and APF do not have. Therefore the comparison
is between a learned controller with map memory and classical
controllers without it, and on equal information APF will win at both
budgets tested. A third result came from auditing the apparatus, i.e.,
repeating training at four seeds will provide a 7.8-point across-seed
spread at one million steps, which will halve by 1.85M, where every
seed will pass APF and the policy will reach 98.8\% at 20\% density.

\textbf{Contributions.} There are five. 1. A four-way comparison on
identical matched layouts with 50 replications, paired significance
tests, effect sizes, and bootstrap intervals, none of which will be
reported by any of the eight studies surveyed in
Table~\ref{tab:related}. 2. A controlled measurement of the
credit-assignment problem, in which the reward scheme represents the
only variable, and which will cost more coverage than the gap between
any two controllers here. 3. Reporting of variance across training
runs, which the surveyed literature omits. 4. Evidence that percentage
coverage over free cells will mislead when free area varies. 5. A
working decentralized MARL coverage system, with the training
configuration that makes it learnable identified explicitly.

\textbf{Limitations.} The simulation represents a 2D kinematic grid
with perfect sensing and localization, and no inertia, wind, or
battery model. Therefore only the relative ordering will plausibly
transfer. In addition, the learned controller will carry memory
features that the classical controllers lack, and the ablation shows
that these will lift its ceiling. Thus the headline comparison is not
one of equal information. The main experiment will also report one
training seed per condition, at a budget that the seed study shows to
be undertrained, and that study will cover 20\% density only.
Therefore the density-dependent claims are not settled. Swarm size was
fixed at six, hyperparameters were tuned only in the rollout-length
dimension, the sensor-noise experiment is withdrawn, and a separate
policy will be trained per density, rather than one generalizing
across them.

\textbf{Future work.} The immediate step will be to repeat the
four-way comparison with several training seeds per cell at a budget
past convergence, which the 20\% result suggests would turn parity
into an advantage for the learned controller. Beyond that, one policy
trained across all densities would test generalization over
environment statistics as well as layouts. In addition, minimal
communication would test whether the 30\% deficit represents a
partial-observability limit or a policy limit, and a valid
sensor-noise experiment would provide the robustness ordering that the
withdrawn one could not. Longer term, PyBullet dynamics and Crazyflie
flights at the Bristol Robotics Laboratory would test whether the
ordering will survive hardware.

%% ===================== BIBLIOGRAPHY (excluded from count) ======
% References set one step down, as IEEE conference papers normally do.
% All 28 entries are unchanged and complete; only the type size moves.
{\footnotesize\linespread{0.90}\selectfont
\bibliographystyle{IEEEtran}
\bibliography{References}
}

\end{document}
